% GENERATED FILE. Edit canonical lessons, then run npm run generate:books.
% canonical-source-hash: fnv1a64:ff9239c3efdd7b9c
% canonical-lessons: GE-C61-angst, GE-C61-freude, GE-C61-sohn, GE-C61-tochter, GE-C61-grossvater

\chapter{Fear, Joy, Son, Daughter, Grandfather}
\label{ch:ge-fear-joy-son-daughter-grandfather}
\hlchaptermodality{61}

\begin{chapteropening}
\textbf{By the end of this chapter:} \emph{I can say fear, joy, a son, a daughter and a grandfather.}

Complete the last lesson of chapter 61: I can say fear, joy, a son, a daughter and a grandfather.
\end{chapteropening}

\section[die Angst]{die Angst — fear}

\label{lesson:GE-C61-angst}



\begin{quote}
Before the new one: say the German for sad, then the German for happy.
\end{quote}

\subsection*{You'll want to know: die Angst}
\textbf{die Angst} — \textquotedblleft{}fear\textquotedblright{}.

\begin{cousinweb}
\textbf{Angst} $\leftarrow$ Old High German \emph{angust} $\leftarrow$ Proto-Germanic \emph{*angustiz}, from \emph{*anguz}, \textquotedblleft{}narrow\textquotedblright{} — inherited: fear as a tightness in the chest. (EtymDB 2.1) English later borrowed \emph{angst} from German.
\end{cousinweb}

\begin{grammarlens}[title={What you've built}]
One more word for this chapter.
\end{grammarlens}

\subsection*{Your turn}
\begin{itemize}
\raggedright
  \item \emph{Say it:} \emph{die Angst}
  \item \emph{Say it:} \emph{die Angst}, once more
  \item \emph{Say it:} say \emph{die Angst}
  \item \emph{Recall:} say the German for sad, then the German for happy, then say \emph{die Angst} again
\end{itemize}

\begin{tcolorbox}[breakable,colback=teal!4,colframe=teal!35!black,title={Before you move on}]
What does die Angst mean? (\textquotedblleft{}Fear\textquotedblright{}.) Say it once more.
\end{tcolorbox}
\section[die Freude]{die Freude — joy}

\label{lesson:GE-C61-freude}



\begin{quote}
Before the new one: say the German for happy, then the German for fear.
\end{quote}

\subsection*{You'll want to know: die Freude}
\textbf{die Freude} — \textquotedblleft{}joy\textquotedblright{}.

\begin{cousinweb}
\textbf{Freude} $\leftarrow$ Old High German \emph{frouwida} (EtymDB 2.1), a noun built on the old adjective for \textquotedblleft{}glad\textquotedblright{}, Old High German \emph{frō} — inherited.
\end{cousinweb}

\begin{grammarlens}[title={What you've built}]
One more word for this chapter.
\end{grammarlens}

\subsection*{Your turn}
\begin{itemize}
\raggedright
  \item \emph{Say it:} \emph{die Freude}
  \item \emph{Say it:} \emph{die Freude}, once more
  \item \emph{Say it:} say \emph{die Freude}
  \item \emph{Recall:} say the German for happy, then the German for fear, then say \emph{die Freude} again
\end{itemize}

\begin{tcolorbox}[breakable,colback=teal!4,colframe=teal!35!black,title={Before you move on}]
What does die Freude mean? (\textquotedblleft{}Joy\textquotedblright{}.) Say it once more.
\end{tcolorbox}
\section[der Sohn]{der Sohn — a son}

\label{lesson:GE-C61-sohn}



\begin{quote}
Before the new one: say the German for fear, then the German for joy.
\end{quote}

\subsection*{You'll want to know: der Sohn}
\textbf{der Sohn} — \textquotedblleft{}a son\textquotedblright{}.

\begin{cousinweb}
\textbf{Sohn} $\leftarrow$ Old High German \emph{sunu} $\leftarrow$ Proto-Germanic \emph{*sunuz} — inherited, the same word as English \emph{son}. (EtymDB 2.1)
\end{cousinweb}

\begin{grammarlens}[title={What you've built}]
One more word for this chapter.
\end{grammarlens}

\subsection*{Your turn}
\begin{itemize}
\raggedright
  \item \emph{Say it:} \emph{der Sohn}
  \item \emph{Say it:} \emph{der Sohn}, once more
  \item \emph{Say it:} say \emph{der Sohn}
  \item \emph{Recall:} say the German for fear, then the German for joy, then say \emph{der Sohn} again
\end{itemize}

\begin{tcolorbox}[breakable,colback=teal!4,colframe=teal!35!black,title={Before you move on}]
What does der Sohn mean? (\textquotedblleft{}A son\textquotedblright{}.) Say it once more.
\end{tcolorbox}
\section[die Tochter]{die Tochter — a daughter}

\label{lesson:GE-C61-tochter}



\begin{quote}
Before the new one: say the German for joy, then the German for a son.
\end{quote}

\subsection*{You'll want to know: die Tochter}
\textbf{die Tochter} — \textquotedblleft{}a daughter\textquotedblright{}.

\begin{cousinweb}
\textbf{Tochter} $\leftarrow$ Old High German \emph{tohter} $\leftarrow$ Proto-Germanic \emph{*duhtēr} — inherited, the same word as English \emph{daughter}. (EtymDB 2.1)
\end{cousinweb}

\begin{grammarlens}[title={What you've built}]
One more word for this chapter.
\end{grammarlens}

\subsection*{Your turn}
\begin{itemize}
\raggedright
  \item \emph{Say it:} \emph{die Tochter}
  \item \emph{Say it:} \emph{die Tochter}, once more
  \item \emph{Say it:} say \emph{die Tochter}
  \item \emph{Recall:} say the German for joy, then the German for a son, then say \emph{die Tochter} again
\end{itemize}

\begin{tcolorbox}[breakable,colback=teal!4,colframe=teal!35!black,title={Before you move on}]
What does die Tochter mean? (\textquotedblleft{}A daughter\textquotedblright{}.) Say it once more.
\end{tcolorbox}
\section[der Großvater]{der Großvater — a grandfather}

\label{lesson:GE-C61-grossvater}



\begin{quote}
Before the new one: say the German for a son, then the German for a daughter.
\end{quote}

\subsection*{You'll want to know: der Großvater}
\textbf{der Großvater} — \textquotedblleft{}a grandfather\textquotedblright{}.

\begin{cousinweb}
\textbf{Großvater} is the word for \textquotedblleft{}great, big\textquotedblright{} + \emph{Vater}: a \textquotedblleft{}great father\textquotedblright{}, modelled on French \emph{grand-père}. The \textquotedblleft{}great\textquotedblright{} word $\leftarrow$ Old High German \emph{grōz} $\leftarrow$ Proto-Germanic \emph{*grautaz}. (EtymDB 2.1)
\end{cousinweb}

\begin{grammarlens}[title={What you've built}]
One more word for this chapter.
\end{grammarlens}

\subsection*{Your turn}
\begin{itemize}
\raggedright
  \item \emph{Say it:} \emph{der Großvater}
  \item \emph{Say it:} \emph{der Großvater}, once more
  \item \emph{Say it:} say \emph{der Großvater}
  \item \emph{Recall:} say the German for a son, then the German for a daughter, then say \emph{der Großvater} again
\end{itemize}

\begin{tcolorbox}[breakable,colback=teal!4,colframe=teal!35!black,title={Before you move on}]
What does der Großvater mean? (\textquotedblleft{}A grandfather\textquotedblright{}.) Say it once more.
\end{tcolorbox}
